\chapter{Schur--Weyl Duality}

Let \(V\) be a complex vector space of dimension \(d\), and put
\(X=V^{\otimes n}\). The two actions on \(X\) exchange the tensor
positions and change the vectors in all positions simultaneously.
Their images in \(\operatorname{End}(X)\) are each other's centralizers.
We first prove the two centralizer identities, then determine the
summands and their characters.

\section{Commuting actions and centralizers}

\begin{proposition}[The two tensor actions]
\label{res:D01}
The symmetric group \(S_n\) and the group \(\operatorname{GL}(V)\)
act linearly on \(X\) by
\begin{align*}
 P_\sigma(v_1\otimes\cdots\otimes v_n)
   &=v_{\sigma^{-1}(1)}\otimes\cdots\otimes v_{\sigma^{-1}(n)},\\
 G_g(v_1\otimes\cdots\otimes v_n)
   &=gv_1\otimes\cdots\otimes gv_n .
\end{align*}
They commute. Thus the images
\[
 A=\operatorname{im}\bigl(\mathbb C[S_n]\longrightarrow
                         \operatorname{End}(X)\bigr),\qquad
 B=\operatorname{span}_{\mathbb C}\{G_g:g\in\operatorname{GL}(V)\}
\]
are commuting subalgebras of \(\operatorname{End}(X)\).
\end{proposition}

\begin{proof}
The displayed expressions are multilinear in the input vectors.
Under successive permutations \(\tau\) and \(\sigma\), the input
at output position \(i\) comes from \(\tau^{-1}\sigma^{-1}(i)\).
This proves \(P_\sigma P_\tau=P_{\sigma\tau}\), and the identity
permutation acts as the identity. The group law
\(G_gG_h=G_{gh}\) holds in every position. The two orders of action
give \(gv_{\sigma^{-1}(i)}\) at position \(i\), so they agree on
the spanning pure tensors. Extending the permutation action
linearly gives the algebra map defining \(A\). Products of the
generators of \(B\) are again generators, and \(G_1\) is its identity.

For \(n=0\) we have \(X=\mathbb C\) with both actions trivial.
For \(d=0<n\) we have \(X=0\), whose endomorphism algebra is the
zero algebra. We use the unique element of that algebra as its
identity. These conventions apply throughout this chapter.
The actions agree with those of \cite[\S\S4.18--4.19]{etingof}.
\end{proof}

\begin{lemma}[Endomorphisms as tensors]
\label{res:D02}
The map
\[
 \Phi\colon\operatorname{End}(V)^{\otimes n}\longrightarrow
                 \operatorname{End}(X),\qquad
 T_1\otimes\cdots\otimes T_n\longmapsto
 \bigl(v_1\otimes\cdots\otimes v_n\mapsto T_1v_1\otimes\cdots\otimes T_nv_n\bigr)
\]
is an algebra isomorphism. It identifies conjugation by \(P_\sigma\)
with permutation of the tensor positions in the source.
Consequently \(\operatorname{End}_{A}(X)\) corresponds to the
\(S_n\)-invariant tensors.
\end{lemma}

\begin{proof}
For \(d,n\geq1\), choose a basis \(v_1,\ldots,v_d\) of \(V\)
and let \(E_{ab}\) send \(v_b\) to \(v_a\) and all other basis
vectors to zero. With \(\mathbf a,\mathbf b\in\{1,\ldots,d\}^n\),
\(E_{a_1b_1}\otimes\cdots\otimes E_{a_nb_n}\) maps to the matrix
unit in the tensor basis taking the word \(\mathbf b\) to
\(\mathbf a\). Thus \(\Phi\) takes a basis to a basis.
Composition of pure tensor operators is multiplication position
by position, and the tensor of the identities is the identity.
Evaluation of \(P_\sigma\Phi(T_1\otimes\cdots\otimes T_n)P_{\sigma^{-1}}\)
on a pure tensor gives \(T_{\sigma^{-1}(i)}\) at position \(i\).
An operator commutes with \(A\) exactly when it commutes with
every \(P_\sigma\), which proves the invariance criterion.
For \(n=0\) the map is the scalar identification, and for \(d=0<n\)
both vector spaces are zero. This is the comparison used in
\cite[Theorem 4.55]{etingof}.
\end{proof}

\begin{theorem}[The first centralizer identity]
\label{res:D03}
The centralizer of \(A\) in \(\operatorname{End}(X)\) is \(B\).
\end{theorem}

\begin{proof}
For \(n=0\) the two algebras and the ambient endomorphism algebra
are \(\mathbb C\), and for \(d=0<n\) all are zero. Suppose \(d,n\geq1\)
and set \(W=\operatorname{End}(V)\). In \(W^{\otimes n}\) every
invariant tensor is its own average over \(S_n\). Since pure
tensors span, the invariant tensors are spanned by sums of
permutations of pure tensors. For \(T_1,\ldots,T_n\in W\),
multilinearity gives
\begin{equation*}
 \sum_{J\subseteq\{1,\ldots,n\}} (-1)^{n-|J|}
      \biggl(\sum_{j\in J}T_j\biggr)^{\otimes n}
  =\sum_{\sigma\in S_n}
         T_{\sigma(1)}\otimes\cdots\otimes T_{\sigma(n)} .
\end{equation*}
Indeed a word in indices using exactly the set \(I\) occurs on the
left with coefficient \(\sum_{J\supseteq I}(-1)^{n-|J|}\).
This coefficient is zero for a proper subset \(I\) and is one for
the full set. A word of length \(n\) using the full set uses each
index once. The invariant tensors are therefore spanned by the
powers \(T^{\otimes n}\), since these powers are invariant.

For any \(T\in W\), the determinant of \(T+tI\) is a monic
polynomial of degree \(d\). There exist \(n+1\) distinct
\(t_0,\ldots,t_n\) for which \(T+t_jI\) is invertible.
The vector-valued polynomial \(t\mapsto(T+tI)^{\otimes n}\)
has degree at most \(n\). The scalar Lagrange interpolation
formula, applied in each coordinate, expresses its value at
zero as a linear combination of its values at the chosen
points. Under \(\Phi\) those values are generators of \(B\).
Lemma~\ref{res:D02} now gives the required centralizer.
Compare \cite[Lemma 4.56 and Proposition 4.58]{etingof}.
\end{proof}

\begin{theorem}[The second centralizer and the tensor decomposition]
\label{res:D04}
Write
\[
 L_\lambda(V)=\operatorname{Hom}_{S_n}(S^\lambda,X),
 \qquad (g\cdot f)(u)=G_g f(u)
 \quad (g\in\operatorname{GL}(V)).
\]
Evaluation is an isomorphism of \(S_n\times\operatorname{GL}(V)\)
representations
\[
 \bigoplus_{\lambda\vdash n} S^\lambda\otimes L_\lambda(V)
        \longrightarrow X .
\]
Every nonzero \(L_\lambda(V)\) is irreducible for
\(\operatorname{GL}(V)\), and these nonzero modules are mutually
nonisomorphic. The centralizer of \(B\) is \(A\).
\end{theorem}

\begin{proof}
By Maschke and Schur \cite[Chapter 8]{algebra-core} and the
classification in \ref{res:A07}, \(X\) is a direct sum of copies
of the \(S^\lambda\). On a copy of \(S^\lambda\), evaluation
from \(S^\lambda\otimes\operatorname{Hom}_{S_n}(S^\lambda,S^\lambda)\)
is scalar evaluation and hence an isomorphism. The Hom spaces
between different simple modules vanish. Applying these facts to
each copy proves the evaluation isomorphism. Postcomposition
defines the asserted action on \(L_\lambda(V)\) because \(G_g\)
commutes with \(S_n\). Evaluation is equivariant for both actions.

Put \(R=\mathbb C[S_n]\). We first identify its full image on
the simple modules. The action map
\[
 R\longrightarrow \prod_{\lambda\vdash n}\operatorname{End}(S^\lambda)
\]
is injective because any element that annihilates every simple
module annihilates the semisimple left regular module and hence
annihilates \(1\) in that module. The multiplicity of \(S^\lambda\)
in the regular module is \(\dim S^\lambda\). Indeed evaluation
at \(1\) identifies \(\operatorname{Hom}_R(R,S^\lambda)\) with
\(S^\lambda\), and computing the same Hom space with the
regular module decomposed into simple summands gives the
multiplicity by Schur. It follows that
\(\dim R=\sum_{\lambda}(\dim S^\lambda)^2\), so the injective
action map is an isomorphism.

Across the summands \(S^\lambda\otimes L_\lambda(V)\) with
\(L_\lambda(V)\ne0\), this identifies \(A\) with the product of the
independent full matrix blocks \(\operatorname{End}(S^\lambda)\).
Each block acts on the first factor and by the identity on the
second. In particular \(A\) contains the projectors to these
summands. An operator commuting with \(A\) preserves every
summand because it commutes with those projectors. In a basis
\((u_i)\) of \(S^\lambda\), commuting with each diagonal matrix
unit makes the operator preserve every \(u_i\otimes L_\lambda(V)\).
Commuting with each off-diagonal matrix unit identifies the
operators on these spaces with one common
\(T\in\operatorname{End}(L_\lambda(V))\). Conversely
\(I\otimes T\) commutes with the entire first block. By
Theorem~\ref{res:D03}, \(B\) is the product of these independent
second-factor blocks.

An operator commuting with \(B\) preserves the same summands
because \(B\) also contains their projectors. Applying the matrix
unit argument to the second factor shows that the new centralizer
is \(A\). Since the linear span of the group on \(L_\lambda(V)\)
is the full matrix algebra, a nonzero invariant subspace contains
the image of each of its vectors under all linear maps, so it is the
whole space. A group intertwiner between two such spaces intertwines
the whole span \(B\), whose projector onto one summand is the
identity there and zero on the other. It must be zero.
If \(X=0\), the evaluation and centralizer identities are between
zero spaces or algebras. The block argument gives the
double centralizer conclusion of
\cite[Theorems 4.54 and 4.57 and Corollary 4.59]{etingof}.
\end{proof}

\section{The length bound and the group algebra kernel}

\begin{theorem}[The exact length bound]
\label{res:D05}
For \(\lambda\vdash n\), \(L_\lambda(V)\) is nonzero if and only
if the number of nonzero parts \(\ell(\lambda)\) is at most \(d\).
\end{theorem}

\begin{proof}
For a set \(C\) of positions let \(b_C\in\mathbb C[S_n]\) be the
sum of the permutations of \(C\), weighted by their signs and
extended by the identity off \(C\). If \(|C|>d\), each basis
tensor in \(X\) has a repeated basis vector among those positions.
The transposition of these two positions fixes that tensor.
Pairing permutations of \(C\) that differ on the right by this
transposition shows that \(b_C\) sends the tensor to zero.
The same conclusion holds if there are no basis tensors because
\(X=0\).

If \(\ell(\lambda)>d\), take \(C\) to be the first column of a
tableau \(t\) of shape \(\lambda\). The polytabloid \(e_t\) in
\ref{res:A04} transforms
by the sign under each permutation of \(C\), so
\(b_C e_t=|C|!\,e_t\). An equivariant map from \(S^\lambda\) to
\(X\) sends this multiple to zero.
The scalar is nonzero, and \(e_t\) generates \(S^\lambda\), so
the map itself is zero.

Suppose \(\ell(\lambda)\leq d\). For a tabloid \(\{t\}\), take the
tensor with the \(i\)-th basis vector of \(V\) at every position whose
label belongs to row \(i\) of \(t\). This construction is unchanged
by reordering entries within a row. Distinct tabloids give
distinct tensor basis words, since the word recovers each row's
label set. Formula \ref{res:D01} for the left action shows that
relabeling the tabloid permutes tensor positions equivariantly.
These words span the subspace of \(X\) of content \(\lambda\)
in the chosen basis, and the tabloid map identifies that subspace
with \(M^\lambda\). It therefore restricts to an injection
of \(S^\lambda\) by \ref{res:A04}. The empty shape for \(n=0\)
gives the empty word, a basis of \(X=\mathbb C\).
The length criterion also appears in
\cite[Theorem 4.63]{etingof}, where it follows from the character.
\end{proof}

\begin{corollary}[The alternating generator of the kernel]
\label{res:D06}
In the full matrix block decomposition of \(\mathbb C[S_n]\),
the kernel of its action on \(X\) is the sum of the blocks with
\(\ell(\lambda)>d\). For \(n>d\) this is the two-sided ideal
generated by the alternating sum \(b_{\{1,\ldots,d+1\}}\).
For \(n\leq d\) the action of the group algebra is injective.
\end{corollary}

\begin{proof}
Theorem~\ref{res:D04} identifies the action with the product of
the blocks with \(L_\lambda(V)\ne0\). Theorem~\ref{res:D05}
identifies those omitted. For \(n>d\), the pairing argument
of Theorem~\ref{res:D05} shows that the indicated alternating
sum acts as zero on \(X\). In each omitted block choose a
tableau \(t\) with \(1,\ldots,d+1\) in its first column.
The alternating sum sends \(e_t\) to \((d+1)!\,e_t\), so its
projection to that block is nonzero. Central block projectors
separate the block components of the two-sided ideal.
A nonzero matrix \(T\) generates the entire two-sided ideal in a
full matrix algebra, since \(E_{ia}TE_{bj}=T_{ab}E_{ij}\) for
a nonzero entry \(T_{ab}\). The ideal is therefore exactly
the sum of all omitted blocks. For \(d=0<n\) the alternating
generator on one position is the unit, as required.
For \(n\leq d\) all partitions of \(n\) have length at most \(d\).
\end{proof}

\section{Characters and dimensions}

\begin{proposition}[The trace along permutation cycles]
\label{res:D07}
For \(g\in\operatorname{GL}(V)\) and \(\sigma\in S_n\),
\[
 \operatorname{tr}_X(P_\sigma G_g)
       =\prod_{c\text{ a cycle of }\sigma}
                   \operatorname{tr}_V(g^{|c|}).
\]
If \(\chi^\lambda\) is the character of \(S^\lambda\), then
\begin{equation}
\label{eq:schur-weyl-trace-average}
 \operatorname{tr}_{L_\lambda(V)}(g)
 =\frac{1}{n!}\sum_{\sigma\in S_n}
     \chi^\lambda(\sigma^{-1})
     \prod_{c\text{ a cycle of }\sigma}\operatorname{tr}_V(g^{|c|}).
\end{equation}
\end{proposition}

\begin{proof}
For \(d,n\geq1\), let \(g_{ab}\) be the matrix entry in output
row \(a\) and input column \(b\). For a basis word with index
\(\mathbf i\), the corresponding diagonal matrix entry of
\(P_\sigma G_g\) is \(\prod_j g_{i_{\sigma(j)},i_j}\).
Summing over the indices divides into independent sums over the
cycles. On a cycle \((j_1,\ldots,j_k)\) the factors are
\(g_{i_{j_2},i_{j_1}}\cdots g_{i_{j_1},i_{j_k}}\), whose sum is
\(\operatorname{tr}_V(g^k)\) by the definition of matrix
multiplication and trace. For \(n=0\) the trace and the empty
product equal one. For \(d=0<n\) they both equal zero.

Using the evaluation decomposition of Theorem~\ref{res:D04},
the same trace is
\(\sum_{\mu\vdash n}\chi^\mu(\sigma)\operatorname{tr}_{L_\mu(V)}(g)\).
Finite-group character orthogonality
\cite[Chapter 8]{algebra-core} gives
\(\frac{1}{n!}\sum_\sigma\chi^\lambda(\sigma^{-1})\chi^\mu(\sigma)
=\delta_{\lambda\mu}\). To see the inverse in this formula,
every \(S_n\)-representation can be made unitary by averaging a
positive definite Hermitian form, so the character on the inverse
is the complex conjugate. Multiplying the decomposition by
\(\chi^\lambda(\sigma^{-1})\) and averaging yields
\eqref{eq:schur-weyl-trace-average}. See \cite[\S4.21]{etingof}.
\end{proof}

\begin{theorem}[Weyl character and dimension]
\label{res:D08}
For invertible \(g=\operatorname{diag}(x_1,\ldots,x_d)\),
\(\operatorname{tr}_{L_\lambda(V)}(g)=s_\lambda(x_1,\ldots,x_d)\),
giving the polynomial character on the diagonal.
If \(\ell(\lambda)\leq d\), pad \(\lambda\) with zeros to length
\(d\). Then
\[
 \dim L_\lambda(V)=
    \prod_{1\leq i<j\leq d}\frac{\lambda_i-\lambda_j+j-i}{j-i}.
\]
For \(\ell(\lambda)>d\) the dimension is zero.
\end{theorem}

\begin{proof}
On the diagonal, \(\operatorname{tr}_V(g^k)=p_k(x)\).
Each permutation is conjugate to its inverse because the two
have the same cycle lengths (\ref{res:A01}). The cycle
class \(\nu\) has \(n!/z_\nu\) elements. Grouping terms in
\eqref{eq:schur-weyl-trace-average} yields
\[
 \operatorname{tr}_{L_\lambda(V)}(g)
    =\sum_{\nu\vdash n}\frac{\chi^\lambda(\nu)}{z_\nu}p_\nu(x)
    =s_\lambda(x)
\]
by the Frobenius characteristic identity \ref{res:C05}. This
also gives the polynomial extension of the diagonal character.
For too many rows the semistandard model \ref{res:C04}
has an empty set of fillings with entries from \(1,\ldots,d\),
in agreement with Theorem~\ref{res:D05}.

For \(d\geq1\) and \(\ell(\lambda)\leq d\), use the
alternant ratio in \ref{res:C04}. Write \(b_j=\lambda_j+d-j\)
and \(a_j=d-j\), both strictly decreasing sequences.
Set \(x_i=t^{d-i}\) for real \(t>1\). The numerator matrix
has entries \((t^{b_j})^{d-i}\), a Vandermonde matrix with rows
indexed by decreasing powers. The denominator is the same
matrix with \(b_j\) replaced by \(a_j\). Thus the ratio is
\(\prod_{i<j}(t^{b_i}-t^{b_j})/(t^{a_i}-t^{a_j})\).
Each factor tends to \((b_i-b_j)/(a_i-a_j)\) as \(t\to1\),
by dividing both differences by \(t-1\). The Schur polynomial
is continuous at \((1,\ldots,1)\), where the trace is the
dimension. This gives the formula. For \(d=0\), the only
permitted shape is empty and its space is \(\mathbb C\).
For the empty shape every factor is one, and for \(d=0,1\)
the product has no factors. These are the formulas of
\cite[Theorem 4.63]{etingof}.
\end{proof}

\begin{example}[Symmetric and exterior powers]
\label{res:D09}
For \(n\geq1\) the row shape and column shape give, respectively,
\[
 L_{(n)}(V)\cong\operatorname{Sym}^n V,\qquad
 L_{(1^n)}(V)\cong\bigwedge^n V.
\]
The zero-degree conventions give \(\mathbb C\) in both cases.
If \(d=2\) and \(n=3\), the tensor decomposition has dimensions
\(8=4+2\cdot2\), with no summand from the shape \((1,1,1)\).
\end{example}

\begin{proof}
By \ref{res:A07} the two simple \(S_n\)-modules are the
trivial and sign representations. Evaluating their intertwining
maps at \(1\) identifies their multiplicity spaces with the
invariant and alternating subspaces, respectively. For
\(\epsilon=1\) or \(\epsilon=\operatorname{sgn}\), set
\[
 Q_\epsilon=\frac{1}{n!}\sum_{\sigma\in S_n}\epsilon(\sigma)P_\sigma .
\]
Reindexing the sum by left or right multiplication gives
\(P_\tau Q_\epsilon=Q_\epsilon P_\tau=\epsilon(\tau)Q_\epsilon\).
Also \(Q_\epsilon\) restricts to the identity on the vectors
transforming by \(\epsilon\). Its kernel is the subspace
spanned by \(v-\epsilon(\sigma)P_\sigma v\), since it kills
these differences and \(v-Q_\epsilon v\) is their average.
Thus it identifies the isotypic subspace with the corresponding
quotient. For \(\epsilon=1\) that quotient is the symmetric
power by the multilinear definition. For the sign it is the exterior
power. An alternating multilinear map kills the signed
permutation differences, by exchanging adjacent positions.
Conversely in the quotient a transposition gives the sign rule and
a tensor with two equal positions is its own negative, hence
zero. These properties give precisely the universal
alternating multilinear map. Both identifications commute
with the diagonal \(\operatorname{GL}(V)\) action.

When \(d=2,n=3\), Theorem~\ref{res:D05} leaves only the
shapes \((3)\) and \((2,1)\). Their Specht dimensions are
\(1,2\) (the standard basis or the hook-length formula
\ref{res:B08}). Theorem~\ref{res:D08} gives multiplicity-space
dimensions \(4,2\). This proves the dimension check and
recovers the symmetric and alternating cases of
\cite[Example 4.60]{etingof}.
\end{proof}

\section{Lie algebra and polynomial representations}

\begin{corollary}[The Lie algebra action]
\label{res:D10}
For \(T\in\mathfrak{gl}(V)=\operatorname{End}(V)\) with the
commutator bracket, let
\[
 \mathcal D(T)=\sum_{i=1}^n
       I^{\otimes(i-1)}\otimes T\otimes I^{\otimes(n-i)}
       \quad\text{acting on }X.
\]
The induced action of its universal enveloping algebra has
image \(B\). In particular its action has the same centralizer
\(A\), and the nonzero multiplicity spaces in
Theorem~\ref{res:D04} are irreducible and mutually
nonisomorphic as \(\mathfrak{gl}(V)\)-modules.
\end{corollary}

\begin{proof}
On distinct positions the summand operators commute.
On a common position their commutator acts by \([T,U]\).
This gives \([\mathcal D(T),\mathcal D(U)]=\mathcal D([T,U])\).
Extending the linear map to the tensor algebra and imposing
the defining commutator relations yields a unital algebra map
from the universal enveloping algebra. Write \(B'\) for its image.

For \(d,n\geq1\), \(\mathcal D(T)\) is the coefficient of \(t\)
in \((I+tT)^{\otimes n}\). Choose \(n+1\) distinct \(t\)'s
where \(I+tT\) is invertible. Interpolation at these points
expresses that coefficient as a linear combination of operators
in \(B\). This gives \(B'\subseteq B\).

Every invertible complex matrix is a product of exponentials.
Gaussian elimination expresses an invertible matrix as a
product of coordinate shears, invertible diagonal matrices
and coordinate swaps. A shear \(I+aE_{ij}\) for \(i\ne j\)
is \(\exp(aE_{ij})\), and an invertible diagonal matrix is
the exponential of the diagonal matrix of chosen scalar
logarithms. On two coordinates a swap matrix is a product of
three shears and a diagonal matrix, as seen from
\[
 \begin{pmatrix}1&1\\0&1\end{pmatrix}
 \begin{pmatrix}1&0\\-1&1\end{pmatrix}
 \begin{pmatrix}1&1\\0&1\end{pmatrix}
 \begin{pmatrix}-1&0\\0&1\end{pmatrix}
     =\begin{pmatrix}0&1\\1&0\end{pmatrix}.
\]
The summands defining \(\mathcal D(T)\) commute, so multiplying
their absolutely convergent exponential series gives
\(\exp(\mathcal D(T))=(\exp T)^{\otimes n}\). The exponential
belongs to \(B'\) because its finite sums belong to the unital
subalgebra \(B'\), and every finite-dimensional linear subspace
is closed. Hence \(G_g\in B'\) for every \(g\in\operatorname{GL}(V)\).
This proves \(B=B'\). For \(n=0\), the derived operators vanish
and the unital image is the scalar algebra, and for \(d=0<n\)
both images are zero. The centralizer and simplicity assertions
now follow from the block argument of Theorem~\ref{res:D04},
also proving the Lie algebra version in
\cite[Theorems 4.55 and 4.57]{etingof}.
\end{proof}

\begin{theorem}[Polynomial representations of the general linear group]
\label{res:D11}
A finite-dimensional representation of \(\operatorname{GL}(V)\)
is called polynomial if its matrix entries in a basis are
polynomials in the matrix entries of \(g\). Every such
representation splits into homogeneous degrees. At any fixed
degree \(n\), the representations are completely reducible with
simple modules exactly \(L_\lambda(V)\) for \(\lambda\vdash n\)
with \(\ell(\lambda)\leq d\).
\end{theorem}

\begin{proof}
First suppose \(d\geq1\) and let \(W\) carry a polynomial
representation \(\rho\). The scalar subgroup gives a finite
expansion \(\rho(tI)=\sum_{k=0}^{N}t^k P_k\) for \(t\ne0\).
Comparing coefficients in the polynomial identity
\(\rho(tI)\rho(sI)=\rho(tsI)\) for \(t,s\ne0\) gives
\(P_iP_j=0\) for distinct \(i,j\) and \(P_i^2=P_i\).
Their sum is \(I_W\) by setting \(t=1\), so their images form a
direct sum. Each \(P_k\) commutes with \(\rho(g)\), as follows
by comparing coefficients in the commutation identity for
scalar matrices. On \(W_k=\operatorname{im}P_k\) we have
\(\rho(tg)|_{W_k}=t^k\rho(g)|_{W_k}\). The matrix entries of the
restricted action are polynomials obtained by a fixed linear
inclusion and projection. Decomposing them by degree and
comparing coefficients of \(t\) shows that only the terms of
degree \(k\) contribute on \(\operatorname{GL}(V)\). This gives
the homogeneous degree decomposition.

Suppose that \(\rho\) is homogeneous of degree \(n\).
The matrix entries of \(G_g\) on the tensor basis are the
products \(\prod_j g_{a_jb_j}\). Every monomial of total degree
\(n\) occurs in this way, including the constant at \(n=0\).
If \(\sum_j c_jG_{g_j}=0\) in \(B\), each of these monomials
has weighted sum zero over the \(g_j\). Every entry of \(\rho\)
is a linear combination of the monomials, so
\(\sum_j c_j\rho(g_j)=0\). Thus \(G_g\mapsto\rho(g)\) extends
to a well-defined linear map from \(B\) to
\(\operatorname{End}(W)\). It respects multiplication because
the group actions do, and sends \(G_1\) to the identity.
This makes \(W\) a \(B\)-module. Conversely every
\(L_\lambda(V)\ne0\) has a homogeneous polynomial action of
degree \(n\): use the evaluation isomorphism to embed it in
\(X\) by tensoring with a fixed nonzero vector of \(S^\lambda\).
Composition with a fixed linear left inverse reads off
its matrix entries from those of \(G_g\).

By Theorems~\ref{res:D04} and \ref{res:D05}, \(B\) is the
product of the full matrix algebras on these spaces. Its
central projectors split every module into its parts for the
single blocks. To see the structure for one block
\(\operatorname{Mat}_{m}(\mathbb C)\), take a basis of the image
of \(E_{11}\) in a module. For each basis vector \(w\), the
vectors \(E_{i1}w\) form a copy of the standard module by the
matrix unit multiplication rule. These vectors for all \(w\)
are independent, as applying \(E_{1i}\) recovers the
coefficients of the \(i\)-th part. They span since
\(\sum_i E_{i1}E_{1i}=1\). Therefore the simple modules of \(B\)
are exactly the standard modules of the occurring blocks and
every module is a direct sum of them. Restricting to the group
gives the required result. For \(d=0\), the general linear
group is trivial and the only polynomials in the empty set of
matrix entries are constants. All nonzero representations
have degree zero, where the unique simple module is
\(L_{\varnothing}(0)=\mathbb C\).
\end{proof}

This is the matrix-entry polynomial case of
\cite[Theorem 4.66(i)]{etingof}. The later algebraic-group topic
treats rational representations (allowing \(\det^{-1}\)),
determinant twists, and algebraic Peter--Weyl
\cite[Theorem 4.66(ii)]{etingof}.
